% Supplementary information for jsac2_v1 (Reflective Intelligent
% Human Bodies). Contents extend the main paper at SI level only --
% no claims live here that aren't summarised in the main text.

\documentclass[10pt,journal,twocolumn]{IEEEtran}

\usepackage{amsmath,amssymb,amsthm,bm}
\usepackage{graphicx}
\usepackage{cite}
\usepackage[T1]{fontenc}
\usepackage{lmodern}
\usepackage{xcolor}
\usepackage{booktabs}
\usepackage{algorithm}
\usepackage{algpseudocode}
\usepackage{hyperref}
\usepackage[capitalise]{cleveref}

\graphicspath{{figures/}{../code/outputs/}}

% Math macros (mirror jsac2_v1.tex).
\newcommand{\rr}{\bm{r}}
\newcommand{\hh}{\mathbf{h}}
\newcommand{\xx}{\mathbf{x}}
\newcommand{\nhat}{\hat{\bm{n}}}
\newcommand{\khat}{\hat{\bm{k}}}
\newcommand{\etahat}{\hat{\bm{\eta}}}
\newcommand{\Sinc}{S_{\mathrm{inc}}}
\newcommand{\Sab}{S_{\mathrm{ab}}}
\newcommand{\Aab}{A_{\mathrm{ab}}}
\newcommand{\Qabs}{\mathbf{Q}_{\mathrm{ab}}}
\newcommand{\Jmat}{\mathbf{J}}
\newcommand{\Phimat}{\mathbf{\Phi}}
\newcommand{\Fn}{\mathbf{F}_n}
\newcommand{\psin}{\bm{\psi}_n}
\newcommand{\thetab}{\bm{\theta}}
\newcommand{\gammab}{\bm{\gamma}}
\newcommand{\hLOS}{h_{\mathrm{LOS}}}
\newcommand{\hbody}{h_{\mathrm{b}}}
\newcommand{\Tepskin}{T_{0}}
\newcommand{\diff}{\mathrm{d}}
\newcommand{\Complex}{\mathbb{C}}
\newcommand{\Reals}{\mathbb{R}}
\newcommand{\ntilde}{\tilde{n}}
\DeclareMathOperator{\RE}{Re}
\DeclareMathOperator{\IM}{Im}

% TAP-style numbering: S1, S2, ... for sections, figures, tables, equations.
\renewcommand{\thesection}{S\arabic{section}}
\renewcommand{\thefigure}{S\arabic{figure}}
\renewcommand{\thetable}{S\arabic{table}}
\renewcommand{\theequation}{S\arabic{equation}}

\theoremstyle{plain}
\newtheorem{proposition}{Proposition}[section]
\theoremstyle{definition}
\newtheorem{definition}[proposition]{Definition}

\title{Supplementary information:\\
Reflective Intelligent Human Bodies}
\author{Robin~Wydaeghe%
  \thanks{R. Wydaeghe is with Ghent University and imec, Ghent, Belgium
  (e-mail: robin.wydaeghe@ugent.be).}}

\begin{document}

\maketitle

% =====================================================================
\section{Approximations 1 and 2: error budgets and proofs}\label{si:approx-proofs}
% =====================================================================

The main paper states two propositions that license the reduction of
the coherent absorption double sum to a single squared norm
(Approximations~1 and~2 of the main paper). The proofs are short and rest on
the same mechanism: the high refractive index of skin compresses all
in-tissue refraction angles to within $\sim\!10^{\circ}$ of normal,
which suppresses both the normal component of the refracted TM field
and the angular dependence of the per-path attenuation and phase
rates. The argument is taken from~\cite{WydaegheMonograph} App.~D and
is reproduced here with the cross-references aligned to the main
paper's symbol set.

\subsection{Proof of Approximation~1 (TM polarisation direction)}

\begin{proposition}[Approximation~1]\label{prop:S-approx1}
Let $\hat{\bm{t}}_n$ be the unit tangent in the incidence plane, so
$\hat{\bm{k}}_n = \sin\theta_n\,\hat{\bm{t}}_n - \cos\theta_n\,\nhat$.
The exact in-tissue TM direction is $\hat{e}'_{p,n} =
\cos\theta_{t,n}\,\hat{\bm{t}}_n + \sin\theta_{t,n}\,\nhat$ with
$\sin\theta_{t,n} = \sin\theta_n / \ntilde$ from Snell's law in
complex form. Replacing $\hat{e}'_{p,n}$ by the incident-side
$\hat{e}_{p,n} = \cos\theta_n\,\hat{\bm{t}}_n + \sin\theta_n\,\nhat$
in the per-path Fresnel operator $\Fn$ leaves self-terms unchanged
and incurs a non-conservative TM--TM cross-term error of
$O(1/|\ntilde|^2)$.
\end{proposition}

\begin{proof}
Self-terms are exact because $\|\hat{e}_{p,n}\|^2 =
\|\hat{e}'_{p,n}\|^2 = 1$. The TM--TM cross-term contains the inner
product
\begin{equation}\label{eq:eppdot}
  \hat{e}'_{p,n} \cdot \hat{e}'_{p,n'}
  = \cos\theta_{t,n}\cos\theta_{t,n'}
  + \sin\theta_{t,n}\sin\theta_{t,n'}.
\end{equation}
Snell's law gives $\sin\theta_{t,n} = \sin\theta_n / \ntilde =
O(1/|\ntilde|)$ across all incidence angles
$\theta_n \in [0^{\circ}, 90^{\circ}]$. The tangential product
$\cos\theta_{t,n}\cos\theta_{t,n'}$ is therefore the dominant
contribution and the normal-normal product is $O(1/|\ntilde|^2)$.
The same expansion applied to the approximated inner product
$\hat{e}_{p,n} \cdot \hat{e}_{p,n'} = \cos\theta_n\cos\theta_{n'}
+ \sin\theta_n\sin\theta_{n'}$ recovers the same dominant term and
substitutes $\sin\theta_n\sin\theta_{n'}$ (which is $O(1)$) for
$\sin\theta_{t,n}\sin\theta_{t,n'}$ (which is $O(1/|\ntilde|^2)$).
The substitution error on the cross-term is therefore
$O(1/|\ntilde|^2)$.

The physical content is that the in-tissue normal field
$E_\perp \propto \sin\theta_t = O(1/|\ntilde|)$, so its contribution
to $\|\bm{E}^{\mathrm{trans}}\|^2$ is $O(1/|\ntilde|^2)$. For skin
at 28~GHz with $|\ntilde| \approx 4.83$, the bound is
$\lesssim 4\,\%$. The error shrinks at higher frequencies and on
tissues with larger $|\ntilde|$;
see~\cite{WydaegheMonograph} Tab.~D.1, reproduced as
Tab.~III (approximation accuracy) of the main paper.
\end{proof}

\subsection{Proof of Approximation~2 (universal depth coupling)}

\begin{proposition}[Approximation~2]\label{prop:S-approx2}
With $\xi_n = \sqrt{\ntilde^2 - 1 + \mu_n^2}$ and the depth
attenuation/phase rates $\alpha_n - j\beta_n = -j k_0 \xi_n$, the
depth coupling
$\Lambda_{nn'} = 1 / [\alpha_n + \alpha_{n'} - j(\beta_{n'} - \beta_n)]$
satisfies
\begin{equation}\label{eq:Gamma-norm}
  \Gamma_{nn'}
    \equiv \frac{\Lambda_{nn'}}{\sqrt{\Lambda_{nn}\Lambda_{n'n'}}}
    = \frac{2\sqrt{\alpha_n \alpha_{n'}}}
           {\alpha_n + \alpha_{n'} - j(\beta_{n'} - \beta_n)},
\end{equation}
with $\Gamma_{nn} = 1$ by construction. For skin at 28~GHz over all
incidence-angle pairs $(\theta_n, \theta_{n'}) \in
[0^{\circ}, 85^{\circ}]^2$, the substitution $\Gamma_{nn'} \approx 1$
incurs a non-conservative error $|1 - \Gamma_{nn'}| \leq 0.44\,\%$
with mean $0.17\,\%$.
\end{proposition}

\begin{proof}
Snell's law in complex form, $\sin\theta_{t,n} = \sin\theta_n/\ntilde$
with $|\ntilde| \approx 4.83$, gives
$|\sin\theta_{t,n}| \leq 1/|\ntilde| \approx 0.21$ across all
incident angles, so $|\theta_{t,n}| \leq \arcsin(0.21) \approx
12^{\circ}$. The tissue refractive index $\ntilde$ therefore
compresses all refracted rays into a $\sim\!10^{\circ}$ cone around
normal. Because $\xi_n = \ntilde \cos\theta_{t,n}$, the variation
of $\xi_n$ across the full incident range is bounded by
$|\xi_n - \ntilde| \leq |\ntilde|(1 - \cos 12^{\circ}) \approx
0.022\,|\ntilde|$. Both $\alpha_n = -k_0 \IM(\xi_n)$ and
$\beta_n = k_0 \RE(\xi_n)$ inherit the same $\sim\!2\,\%$ relative
spread; differences $\alpha_n - \alpha_{n'}$ and $\beta_n - \beta_{n'}$
are second-order in this spread. Substituting into~\eqref{eq:Gamma-norm}
and expanding to first order in
$(\alpha_n - \alpha_{n'})/\bar\alpha$ and
$(\beta_n - \beta_{n'})/\bar\alpha$ (with
$\bar\alpha = \tfrac{1}{2}(\alpha_n + \alpha_{n'})$) gives
\begin{equation}
  |1 - \Gamma_{nn'}|
    \leq \frac{1}{2}
    \left(\frac{\alpha_n - \alpha_{n'}}{\bar\alpha}\right)^2
    + \frac{1}{4}
    \left(\frac{\beta_n - \beta_{n'}}{\bar\alpha}\right)^2
    + O(\delta^3),
\end{equation}
with $\delta \sim 0.02$. The numerical sweep
across $(\theta_n, \theta_{n'}) \in [0^{\circ}, 85^{\circ}]^2$ for
skin at 28~GHz gives
$|1 - \Gamma_{nn'}| \leq 0.44\,\%$ with mean $0.17\,\%$
(see~\cite{WydaegheMonograph} Fig.~D.2). The bound enters the
absorption double sum non-conservatively (it can both over- and
under-estimate cross-term magnitudes) and is an order of magnitude
smaller than the calibration headroom that the body-side $\gammab$
regression in Section~VI (Closed-loop calibration) of the main paper absorbs.
\end{proof}

% =====================================================================
\section{Per-path Fresnel transmission operator $\mathbf{F}_n$}\label{si:Fn}
% =====================================================================

The main paper invokes $\Fn$ as ``rank-2, polarisation-aware''
without writing the matrix explicitly. We collect the formal
definition here. The full Maxwell-from-the-ground-up derivation is
in~\cite{WydaegheMonograph} App.~F (Snell's law from phase matching,
TE and TM boundary conditions on a lossy half-space).

At the surface point $\rr$ on triangle $t$, path $n$ has incidence
cosine $\mu_n = \nhat_t \cdot (-\khat_n)$. For $\mu_n > 0$ the local
TE and TM unit vectors are
\begin{equation}\label{eq:S-TE-TM-basis}
  \hat{e}_{s,n} = \frac{\khat_n \times \nhat_t}
                       {\|\khat_n \times \nhat_t\|},
  \qquad
  \hat{e}_{p,n} = \hat{e}_{s,n} \times \khat_n.
\end{equation}
The Fresnel TE and TM amplitude transmission coefficients on the
lossy half-space, derived in~\cite{WydaegheMonograph} App.~F via
continuity of the tangential $\bm{E}$ and $\bm{H}$ fields at $z=0$,
are
\begin{equation}\label{eq:S-fresnel-t}
  t_{s,n} = \frac{2 \mu_n}{\mu_n + \xi_n},
  \qquad
  t_{p,n} = \frac{2\,\ntilde\,\mu_n}{\ntilde^2\,\mu_n + \xi_n},
\end{equation}
with $\xi_n = \sqrt{\ntilde^2 - 1 + \mu_n^2}$ and $\RE(\xi_n) > 0$.

\begin{definition}[Fresnel transmission operator]\label{def:S-Fn}
The dimensionless rank-2 operator
\begin{equation}\label{eq:S-Fn-def}
  \Fn(\rr) =
  \bigl(t_{s,n}\,\hat{e}_{s,n}\hat{e}_{s,n}^{\!\top}
      + t_{p,n}\,\hat{e}_{p,n}\hat{e}_{p,n}^{\!\top}\bigr)\,H(\mu_n)
\end{equation}
acts on the incident polarisation amplitude vector
$\psin = \psi_{s,n}\hat{e}_{s,n} + \psi_{p,n}\hat{e}_{p,n}$ by
projecting onto the local TE and TM directions and scaling each
component by its Fresnel coefficient. The Heaviside factor
$H(\mu_n)$ vanishes on back-facing paths.
\end{definition}

The operator is rank-2 because the incident-side polarisation basis
$(\hat{e}_{s,n}, \hat{e}_{p,n})$ spans the 2-D transverse plane to
$\khat_n$. In that basis $\Fn$ is diagonal:
\begin{equation}\label{eq:S-Fn-diag}
  \Fn = \mathrm{diag}(t_{s,n}, t_{p,n}) \cdot H(\mu_n),
\end{equation}
where the diagonal form follows from Approximation~1
(\cref{prop:S-approx1}): replacing the refracted TM direction by the
incident TM direction makes the TE and TM channels orthogonal in the
$(\hat{e}_{s,n}, \hat{e}_{p,n})$ basis. The transmitted surface
field is then
\begin{equation}\label{eq:S-Etrans}
  \bm{E}_n^{\mathrm{trans}}(\rr, 0)
    = \Fn(\rr)\,\psin\,x_{j(n)}\,e^{-jk_0\khat_n\cdot\rr},
\end{equation}
which is the per-triangle integrand consumed by the Kirchhoff render
of Section~IV (UE-anchored Kirchhoff render) and by the exposure
operator $\Qabs(\thetab) = \Jmat^{H}\mathbf{M}(\thetab)\Jmat$ of
the main paper.

% =====================================================================
\section{Scene and signal-budget conventions}\label{si:scene}
% =====================================================================

We collect here the scene parameters that the main paper invokes
without re-deriving.

\textbf{Coordinates.} World frame is right-handed Z-up. The user
stands at $\bm{r}_u = (30, 0, 1.2)\,$m. The base station occupies
one of three regimes: far-BS at $\bm{r}_\mathrm{BS} = (0, 0, 8)\,$m
(30~m horizontal, 8~m mast), mid-BS at $(20, 0, 5)\,$m, close-BS
at $(25, 0, 3)\,$m. The phone is held at chest height,
$40\,$cm in front of the body centroid, at world position
$\bm{r}_p = \bm{r}_u + (-0.40, 0, 0.05)\,$m so that its broadside
faces the BS in all three regimes.

\textbf{Body orientation.} The SMPL-X canonical mesh
(Y-up, body facing $+z$) is mapped to the JSAC2 world frame by
$\mathbf{R}_z(-90^{\circ}) \mathbf{R}_x(-90^{\circ})$. This places
the body's vertical axis along world $+z$, the body's facing
direction along world $-x$ (toward the BS), the body's left arm
along world $-y$ and the body's right arm along world $+y$.

\textbf{Link budget.} BS transmit power $P_\mathrm{tx} = 43\,$dBm,
noise floor $-94\,$dBm (corresponding to $100\,$MHz BW with $5\,$dB
noise figure), carrier $f_c = 28\,$GHz, bandwidth $100\,$MHz. The
$8 \times 8$ URA at the BS uses $\lambda/2$ spacing. The main paper's
numerical study (Section~VIII of the main text) is built on
the canonical Sionna RT Munich scene with $16$ traced Rx positions.
This SI documents an additional \emph{parametric scalar-loss
companion} that the main paper invokes for the comfort-vs-rate Pareto
and the baselines comparison: a free-space LOS link with three
scalar-loss bands enumerating the rest of the propagation budget,
$35\,$dB (LOS-attenuated, body is a small perturbation),
$55\,$dB (body-comparable, body-mediated and LOS contributions
comparable in magnitude), and $80\,$dB (body-dominant, body channel
dominates direct LOS). The companion is a controlled abstraction:
it removes the per-scene RT variance to isolate how rate, comfort
budget, and baseline-precoder choice trade off as a single propagation
budget knob. Rates are reported as plain Shannon
$R = \log_2(1 + \mathrm{SINR})$ in bps/Hz throughout.

% =====================================================================
\section{VPoser latent-space details}\label{si:vposer}
% =====================================================================

The variational pose prior of \cite{Pavlakos2019VPoser} parametrises
a learned manifold of plausible body poses. The decoder
$\bm{G}: \Reals^{32} \to \Reals^{63}$ is a 3-layer MLP
($512 \to 512 \to 126$, leaky-ReLU) with continuous-rotation 6-D
output (Zhou et al., 2019) decoded to per-joint axis-angle for the
21 SMPL-X body joints (excluding global orientation and hands).
The encoder $\bm{G}^{-1}$ is a parallel 3-layer MLP returning the
mean and log-variance of an approximate Gaussian posterior; we use
the mean for deterministic AMASS-frame inversion.

\textbf{Pretrained weights.} We load the official V02\_05 checkpoint
(epoch 13) trained on the AMASS corpus~\cite{AMASS2019} with batch
size 128 on the CMU + BMLrub training split.

\textbf{Comfort radius.} The latent ball
$\|\bm{z} - \bm{z}_0\| \leq \rho_z$ with $\rho_z = 1.5$ corresponds,
under $p(\bm{z}) = \mathcal{N}(\bm{0}, \bm{I})$, to a prior
log-probability above $\log p(\bm{z}_0) - 1.125$. Empirically this
threshold rules out the contortionist tail while allowing torso
yaws of $\pm 30^{\circ}$ and shoulder rolls of $\pm 25^{\circ}$ on
the seated/standing baselines we evaluate.

\textbf{Encoder round-trip.} An AMASS walk frame projected to its
posterior mean and decoded back has per-joint reconstruction error
$\sim 11^{\circ}$ on average. This is consistent with the
information-bottleneck loss the VAE was trained against and is
absorbed entirely by the body-side $\gammab$ calibration of the
main paper's Section~V.

% =====================================================================
\section{Calibration recovery}\label{si:calibration}
% =====================================================================

\textbf{Tier-D scalar pre-conditioning.} Before the body-side
$K$-mode SVD fit, we fit a single complex scalar $\beta_0$ on the
LOS coefficient and a single real scalar on $\hbody$ to absorb
gain miscalibration. The Tier-D residual $\bm{r}^{(D)}$ is then
the residual on which the $K$-mode body-side SVD fit operates.

\textbf{Tier-B operator-SVD basis.} The $K = K_{99}$ leading right
singular vectors of $\Lambda_\mathrm{KH}$ form the basis
$\bm{V}_K$. The empirical operator-spectrum sweep of
\cref{si:svd} finds $K_{99} \in \{1, 2, 4\}$ across the three BS
regimes, and the leading mode absorbs $\geq 93\,\%$ of the
body-channel energy at far-BS.

\textbf{Recovery vs.\ SNR sweep on the predecessor's plaza.}
\Cref{fig:resid-snr} reports the post-calibration residual
$\rho = \|\hat{\bm{\beta}} - \bm{\gamma}\| / \|\bm{\gamma}\|$ as a
function of UL pilot SNR for path-set sizes
$N \in \{20, 40, 60\}$ on the predecessor manuscript's plaza
geometry, evaluated under both plain LS and ridge LS
($\lambda = 10^{-2}$). The no-calibration baseline sits at
$-4.7\,$dB (trusting the in-silico amplitudes leaves a
$58\,\%$ relative error on the per-path complex amplitudes). Plain
LS at $N = 20$ beats the no-cal baseline above $15\,$dB UL SNR and
reaches $-21\,$dB at $30\,$dB UL SNR. At $N = 60$, where the column
count approaches the array size ($M = 64$), the steering matrix
develops small singular values that amplify pilot noise; ridge LS
buys back $\sim\!25\,$dB of headroom and beats the no-cal baseline
above $25\,$dB UL SNR. Production calibration should pick
$N < M_\mathrm{ant} / 2$, regularise, or merge correlated paths into
clusters.

\begin{figure}[!t]
\centering
\includegraphics[width=\columnwidth]{residual_vs_snr.pdf}
\caption{Post-calibration residual $\rho$ vs.\ UL pilot SNR on the
predecessor's plaza geometry, $200$ Monte Carlo trials per
$(\mathrm{SNR}, N)$ cell. Solid: plain LS. Dashed: ridge LS,
$\lambda = 10^{-2}$. Dotted horizontal: no-calibration baseline at
$-4.7\,$dB. The plain-LS failure at $N = 60$ comes from
column-correlated steering matrices; ridge regularisation recovers
$\sim\!25\,$dB.}
\label{fig:resid-snr}
\end{figure}

\textbf{Recovery vs.\ scene loss.} At $20\,$dB UL pilot SNR, the
Tier-B fit recovers the body channel to within $1.06\,\%$ of the
ground truth in cap-violation rate, compared to $2.47\,\%$ for the
uncalibrated dual-ascent baseline of \cite{Wydaeghe2026jsacv5}.
The recovery is invariant to scene-loss in the regime sweep above.

% =====================================================================
\section{IMU sensitivity overlay}\label{si:imu}
% =====================================================================

\Cref{fig:imu-overlay} re-runs the prior dual-ascent
exposure-constrained beamforming sensitivity test of
\cite{Wydaeghe2026jsacv5} as a function of the per-joint IMU
attitude noise $\sigma_\mathrm{joint}$. The $\gammab$-corrected
closed loop is overlaid.

\begin{figure}[!t]
\centering
\includegraphics[width=\columnwidth]{imu_gamma_overlay.pdf}
\caption{IMU sensitivity sweep. Cap-violation rate vs.\ per-joint
IMU attitude error $\sigma_\mathrm{joint}$ for four methods. The
$\gammab$-corrected closed loop (green) absorbs $93.5\,\%$ of the
pose-induced channel-prediction error at $\sigma_\mathrm{joint} =
4^{\circ}$ (canonical operating point for an in-pocket
consumer-grade IMU).}
\label{fig:imu-overlay}
\end{figure}

% =====================================================================
\section{Operator spectrum and $K_{99}$ adaptivity}\label{si:svd}
% =====================================================================

\Cref{fig:rank-cdf} reports the $\epsilon$-effective rank of
$\Qabs$ across $20$ sample positions on the Thelonious phantom for
two channel models: a plaza specular dictionary (LOS plus a ground
bounce plus two facade reflections, $4$ paths per body) and a
3GPP~38.901 UMa-LOS stochastic dictionary ($12$ clusters
$\times$ $5$ subpaths $= 56$ subpaths per body), both at $26\,$GHz
on an $8\times 8$ half-wavelength UPA.

At $\epsilon = 10^{-2}$, the median rank is $4$ for the specular
dictionary (max $4$, never above) and $3$ for the stochastic
dictionary (max $3$). The specular case shows a hard cliff after
$k = 4$ to the double-precision numerical floor, which follows from
$\mathrm{rank}(\Qabs) \leq B$ with $B = 4$ bounces. The stochastic
case shows a smooth decay limited by the $\sim\!14^{\circ}$ angular
resolution of the panel: even with $56$ subpaths only $3$ to $5$
distinguishable BS-side directions are excited above $\epsilon =
10^{-2}$, with a long sub-dominant tail that reaches the numerical
floor near $k = 30$.

\Cref{fig:svd-spectrum} reports the singular-value spectrum of
$\Lambda_\mathrm{KH}$ across the three BS regimes of the JSAC2
Munich scene. The far-BS geometry concentrates body-channel energy
in a single mode ($K_{99} = 1$); the close-BS geometry, where the
body subtends a wide solid angle, distributes energy across
$K_{99} = 4$ modes. The two figures cover the same physical
regime---rank-few BS-side excitation of the body---on different
array conventions; the rank-few claim transfers.

\begin{figure}[!t]
\centering
\includegraphics[width=\columnwidth]{rank_cdf.pdf}
\caption{$\epsilon$-effective rank of $\Qabs$ on $20$ Thelonious body
positions at $26\,$GHz with an $8\times 8$ half-wavelength UPA.
Median (blue), per-body traces (grey), $10$--$90$ percentile band
(shaded). Left: plaza specular dictionary, hard cliff at $k = 4$
bounces. Right: 3GPP UMa-LOS stochastic dictionary, smooth decay
limited by array angular resolution. Reproduced from the predecessor
manuscript's rank-check experiment.}
\label{fig:rank-cdf}
\end{figure}

\begin{figure}[!t]
\centering
\includegraphics[width=\columnwidth]{fig2_svd_spectrum.pdf}
\caption{Operator $\Lambda_\mathrm{KH}$ singular-value spectrum
on the JSAC2 Munich scene, normalised per regime. Far-BS is
rank-1; close-BS is rank-4.}
\label{fig:svd-spectrum}
\end{figure}

% =====================================================================
\section{Strict visibility variant for fingertip features}\label{si:visibility}
% =====================================================================

The default Kirchhoff render of Section~IV (UE-anchored Kirchhoff render) of the main
paper culls non-illuminated triangles by a per-path back-face test
$\mu_n = \nhat_t \cdot (-\khat_n) > 0$. The cost is
$O(M \cdot M_\mathrm{tri,vis})$ per pose, with $M_\mathrm{tri,vis}$
the number of front-facing triangles. The back-face test is exact
for convex regions and a faithful approximation for the torso, head,
and limbs at FR2 mesh resolution; it under-counts self-shadowing on
the body's concave features.

\textbf{Strict variant.} For a triangle $t$ to be BS-visible along
path $n$, no other body triangle may block the BS-to-$\bm{c}_t$ ray
segment. The strict visibility flag
$V^{(\mathrm{strict})}_{nt} \in \{0, 1\}$ is the conjunction of the
back-face test with a ray-mesh intersection query:
\begin{equation}\label{eq:V-strict}
  V^{(\mathrm{strict})}_{nt} = H(\mu_n)
    \cdot \bigl[\,\text{no $t' \neq t$ intersects the segment}
                  \bm{c}_t \to \bm{r}_\mathrm{BS}^{(n)}\,\bigr].
\end{equation}
The intersection query is run against all $M_\mathrm{tri,total} =
20\,908$ body triangles per (path, candidate-triangle) pair, so the
per-pose cost scales as
$O(M \cdot M_\mathrm{tri,vis} \cdot M_\mathrm{tri,total})$.

\textbf{Cost factor.} The ratio
$M_\mathrm{tri,total} / M_\mathrm{tri,vis} \approx 5$ on the SMPL-X
mesh ($T = 20\,908$ triangles, with roughly $T_\mathrm{vis}
\approx 4\,000$ front-facing per BS direction). The strict variant
therefore costs $\sim\!5\times$ more per pose than the back-face
default. A bounding-volume hierarchy reduces the per-query
intersection cost to $O(\log M_\mathrm{tri,total})$ but does not
change the asymptotic ratio.

\textbf{When required.} Strict visibility matters when the
self-shadowed feature is comparable in surface area to the
back-face-culled estimate of its visible region. For fingertip,
ear-edge, and inner-thigh features, adjacent body geometry casts a
shadow that the back-face test misses. For torso, head, and
forearm-scale features, the strict and default tests agree to within
the per-triangle sampling resolution. The main paper's numerical
study uses the back-face default; the strict variant is invoked
only for the fingertip-features case.

% =====================================================================
\section{Layered Fabry-Pérot extension to sub-6\,GHz}\label{si:fp}
% =====================================================================

The pseudo-Brewster collapse $\Tepskin \approx 0.54$ at 28~GHz that
licenses the single-Fresnel reduction of the main paper does not
hold below $\sim 10\,$GHz, where layered tissue resonances
(skin/fat/muscle Fabry-Pérot dips) introduce oscillations of
$\pm 30\,\%$ in the per-element transmission coefficient. The
extension recipe is documented in \cite{WydaegheMonograph} Part~II:
replace the constant $\Tepskin$ by the flux-averaged transmission
$\bar{T}(f) = 2\int_0^1 T_\mathrm{avg}(\mu, f)\,\mu\,\diff\mu$
computed once per frequency from a layered-tissue Fresnel sweep.
The body-side $\gammab$ calibration absorbs any residual layered
mismatch at no additional pilot cost.

% =====================================================================
\section{Failure-mode taxonomy}\label{si:failure-modes}
% =====================================================================

The closed loop of Section~VII (Pose-differentiable rate and HITL control) of the main paper either
reaches a comfort-bounded latent perturbation that achieves the
target rate gain or it does not. Failures fall into four categories,
distinguishable from the iterate trace alone.
\Cref{tab:failure-modes} lists each category, the diagnostic
criterion on the iterates $\bm{z}_0, \bm{z}_1, \ldots, \bm{z}_K$,
the typical underlying cause, and the corresponding mitigation.

\begin{table*}[!t]
\centering
\caption{Failure-mode taxonomy for the comfort-bounded latent
gradient ascent. ``BF'' denotes the brute-force optimum within the
comfort ball $\|\bm{z} - \bm{z}_0\| \leq \rho_z = 1.5$, found by
random sampling. ``Best-iterate'' refers to the iterate
$\bm{z}_{k^\star}$ with the highest measured rate, which the closed
loop returns rather than $\bm{z}_K$.}
\label{tab:failure-modes}
\footnotesize
\begin{tabular}{@{}llll@{}}
\toprule
Mode & Diagnostic on iterates & Typical cause & Mitigation \\
\midrule
Stagnant
  & $\|\nabla_{\bm{z}} R\| < \epsilon$ for all $k = 0, \ldots, K$
  & Latent ball is flat at $\bm{z}_0$ (local plateau)
  & Random restart inside comfort ball \\[2pt]
Partial
  & $R(\bm{z}_K) \in [R_\mathrm{BF} - 3, R_\mathrm{BF} - 1]\,$dB
  & FD gradient direction biased by probe noise
  & Larger FD step; more probe averages \\[2pt]
Drift
  & $R(\bm{z}_K) > R(\bm{z}_0)$ but $R(\bm{z}_K) < R(\bm{z}_{k^\star})$
  & Line-search backoff failed to track best-iterate
  & Return best-iterate $\bm{z}_{k^\star}$ \\[2pt]
IMU-floor
  & $R(\bm{z}_K) \approx R(\bm{z}_0)$ at high $\sigma_\mathrm{joint}$
  & FD probe drowns in pose-induced channel-prediction error
  & Require $\sigma_\mathrm{joint} \lesssim 4^{\circ}$ \\
\bottomrule
\end{tabular}
\end{table*}

The mitigations are independent: random restarts address stagnant
ascent without affecting the other modes, larger FD steps cure
partial recovery without affecting drift, best-iterate tracking
fixes drift without affecting IMU-floor, and the IMU-floor mode
sets the operating envelope on $\sigma_\mathrm{joint}$ that the
hardware loop must respect. The empirical study in
Section~VIII of the main paper labels each non-reaching
trajectory by its dominant mode and reports the relative
frequencies stratified by BS distance regime.

% =====================================================================
\section{Compliance metrics}\label{si:compliance}
% =====================================================================

\textbf{Whole-body SAR.} ICNIRP 2020 general-public limit:
$0.08\,$W/kg averaged over $30\,$min over $10\,$g of contiguous
tissue. At our canonical operating point ($43\,$dBm, $30\,$m
range), the whole-body SAR is $\leq 1.2 \times 10^{-4}\,$W/kg,
six orders of magnitude below the limit.

\textbf{Peak local APD.} ICNIRP 2020 ${\geq}\,6\,$GHz general-public
limit: $4\,$W/m$^{2}$ averaged over $4\,$cm$^{2}$. The peak-local
APD over the body surface is $0.58\,\%$ of the limit at the
canonical operating point. The local averaging kernel is a disk
of radius $\sqrt{4\,\mathrm{cm}^{2}/\pi} \approx 1.13\,$cm centred
at each surface vertex; the peak across the body is reported.

\textbf{Why compliance is not the binding constraint.} At FR2
operating points, the link budget (BS-to-body path loss) limits
$\xx$ before any compliance constraint binds. The exposure reading
of the main paper's Section~VIII is therefore informational: it
converts a slack constraint into actionable user-facing telemetry.

\bibliographystyle{IEEEtran}
\begin{thebibliography}{9}
\bibitem{Pavlakos2019VPoser} G.~Pavlakos \emph{et al.}, ``VPoser:
variational human pose prior for body inverse kinematics,''
companion software release, MPI-IS, 2019.
\bibitem{AMASS2019} N.~Mahmood \emph{et al.}, ``AMASS: archive of
motion capture as surface shapes,'' in \emph{Proc.\ ICCV}, 2019.
\bibitem{Wydaeghe2026jsacv5} R.~Wydaeghe, ``Body digital twins for
compliance-aware mmWave downlink,'' technical report (predecessor
manuscript), 2026.
\bibitem{WydaegheMonograph} R.~Wydaeghe, ``A monograph on absorbed
power density on human bodies in wireless environments,'' technical
report, 2026.
\end{thebibliography}

\end{document}
